% introduction.tex — Full proposal draft, Week 8
% Revised from bigpicture.tex (Week 5). Light edits for flow and RQ placement.

\section{Introduction}

Attention is unevenly distributed across populations, and so is our ability to study it. Individual differences in sustained attention and working memory predict a wide range of consequential life outcomes, from academic achievement and occupational performance to vulnerability to clinical conditions like attention-deficit/hyperactivity disorder \citep{kofler2013}. The neuroscientific tools most capable of measuring these differences, however, are expensive, stationary, and concentrated in well-funded research institutions. Functional magnetic resonance imaging (fMRI) scanners cost millions of dollars to purchase and maintain, require specially shielded rooms, and are largely absent from the low-income communities and under-resourced clinical settings where attentional difficulties are most prevalent. As a result, the knowledge base about attention and the brain reflects a structurally skewed sample of who can access a scanner \citep{farah2006}.

This project asks whether that constraint can be relaxed. The central question is whether, if a participant wears a portable brain sensor while watching a movie, the recording from that single naturalistic session can be used to train a computational model that later predicts whole-brain neural activity during structured attention tasks. The portable sensor in question is functional near-infrared spectroscopy (fNIRS), a lightweight headset that measures cortical hemodynamics by detecting changes in light absorption through the scalp. If the approach works, it would allow researchers to approximate fMRI-quality, whole-brain measures of cognitive engagement from a low-cost recording taken in a classroom, a clinic, or a community center, without any scanner access.

The design draws on two datasets linked by a shared naturalistic stimulus, the 1959 Hitchcock film \textit{North by Northwest} (NNW). An ongoing in-lab fNIRS study (N\,=\,28 participants with usable data, spanning participant IDs p01--p34, target N\,=\,40) records oxygenated hemoglobin (HbO) time series from 42 prefrontal, lateral parietal, and temporal-occipital channels while participants watch the film and then complete a series of attention tasks. These tasks include a face-based working memory n-back task at 1-back and 2-back difficulty levels and two rounds of a Continuous Performance Task (CPT). A separate, open fMRI dataset from a collaborating lab (N\,=\,59 different participants) provides fully preprocessed BOLD time series across 122 brain regions from participants who watched the same film in a scanner. Because both groups watched the same movie, their neural responses are temporally aligned even though they are entirely different people. This shared temporal structure is the engine of the approach.

The computational method at the core of the project is adapted principal component regression (aPCR), following \citet{gao2025}. The model first reduces both the fNIRS and fMRI recordings to their most important underlying components via PCA, retaining components capturing 75\% of fNIRS variance (approximately 14 components) and 85\% of fMRI variance (approximately 52 components), matching the thresholds used in the original Gao et al.\ framework. It then uses ordinary least-squares regression to learn the linear mapping between those component spaces. Training uses leave-one-out cross-validation across N\,=\,28 participants. In each fold, 27 fNIRS participants are paired with all 59 fMRI participants (1,593 training pairs), and the held-out fNIRS participant's data are passed through the frozen learned mapping to generate a predicted whole-brain fMRI time series for 122 regions over 886 seconds. The NNW pipeline has been run end-to-end on the full sample. The model significantly predicts 56 of 122 brain regions (FDR-corrected, $q < 0.05$), with particularly strong performance in the dorsal attention network (11/14 ROIs, 79\%) and ventral attention network (14/24 ROIs, 58\%). At the connectivity level, the Pearson correlation between predicted and true inter-subject functional connectivity (ISFC) matrices is $r = 0.617$ ($p = 0.001$, permutation test), which indicates that the model preserves not just local activation patterns but the large-scale organization of brain-wide communication during the film.

The project goes beyond the original \citet{gao2025} framework by reframing the primary question. Rather than asking whether the model can transfer across paradigms in a technical sense, this project poses a social science question: can we detect meaningful individual differences in who pays attention, and can we do so without a scanner? For each cognitive task, the frozen NNW model will be applied to each participant's task fNIRS data to generate predicted whole-brain time series. These predictions will be evaluated for neural plausibility, specifically whether predicted activation patterns in frontoparietal networks exceed default mode network activation during high-load conditions, as the fMRI literature would predict. They will also be evaluated for individual differences, specifically whether per-subject predicted activation patterns correlate with behavioral performance including n-back accuracy, CPT hit rate, and reaction time variability. This is where the equity argument becomes concrete. If a brief, portable, naturalistic recording session can approximate the neural signature of attentional capacity, then rigorous cognitive neuroscience becomes possible in the settings and populations that have historically been excluded from it. Reaction time variability is particularly well-suited as the individual-differences target because it is sensitive enough to discriminate clinical and subclinical attentional profiles, and if the predicted neural patterns track it, the approach offers a path toward accessible, whole-brain-equivalent measures of attentional stability in schools, clinics, and communities where MRI infrastructure simply does not exist.

The project addresses four specific research questions:

\begin{enumerate}
    \item[\textbf{RQ1:}] Does an fNIRS-to-fMRI mapping trained on naturalistic movie-watching generalize to predict whole-brain neural dynamics during structured cognitive tasks (n-back, CPT)?
    \item[\textbf{RQ2:}] Which brain regions maintain reliable fNIRS-to-fMRI mapping across naturalistic and cognitive paradigms, and do these correspond to domain-general attention networks?
    \item[\textbf{RQ3:}] Can the model predict meaningful differences in brain activity between easy and hard working memory conditions, and do those predictions correlate with behavioral performance (accuracy, reaction time variability)?
    \item[\textbf{RQ4:}] Could a single brief movie-watching session serve as a universal calibration paradigm for portable fNIRS, enabling whole-brain-equivalent neurocognitive assessment without a scanner?
\end{enumerate}

This project makes four contributions. First, it provides the first systematic test of whether an fNIRS-to-fMRI mapping trained on naturalistic stimuli generalizes to structured cognitive paradigms without retraining. Second, it introduces a cross-dataset, cross-paradigm validation design that combines a fully processed open fMRI dataset with ongoing portable-sensor data collection, demonstrating a practical template for low-cost neuroimaging research. Third, it targets individual differences in attentional stability as the outcome of interest, grounding the technical test in a consequential social question about who can be measured and where. Fourth, it establishes a concrete equity argument: if the approach succeeds, meaningful whole-brain cognitive neuroscience becomes achievable in schools, clinics, and community settings without MRI infrastructure.
